El circuito de amperímetro, al medir una corriente esperada de $2\,\text{mA}$ se obtuvo una lectura de $1.9928\,\text{mA}$, lo que equivale a un error del $0.36\%$. Este porcentaje tan bajo confirma que la resistencia interna relativa del amperímetro a la del circuito de estudio altera muy poco a este.

En el caso del voltímetro al medir un voltaje de $3\,\text{V}$ en un circuito con resistencias altas ($100\,\text{k}\Omega$), la lectura fue de $2.40\,\text{V}$ en la escala de $12\,\text{V}$ ($20\%$ de error) y de $1.714\,\text{V}$ en la escala de $4\,\text{V}$ ($42.9\%$ de error). En cambio, al reducir el valor de las resistencias del circuito ($10\,\text{k}\Omega$), los valores obtenidos mejoraron a $2.927\,\text{V}$ ($2.43\%$ de error) y $2.791\,\text{V}$ ($6.97\%$ de error). Esto demuestra que mide con mayor exactitud cuando se trabaja con resistencias grandes, ya que cuenta con una resistencia interna mayor.

Finalmente, la simulación en \textit{CircuitJS} permitió comprobar las fórmulas teóricas y la interacción entre el instrumento de medida y el circuito en condiciones ideales, sin los errores asociados a los componentes o las equivocaciones al tomar lecturas en un equipo físico de laboratorio.
